% Generated draft: requires mathematical and source review.
\documentclass[11pt]{article}
\usepackage[margin=25mm]{geometry}
\usepackage{fontspec}
\setmainfont{Latin Modern Roman}
\usepackage{amsmath,amssymb,mathtools,unicode-math}
\setmathfont{Latin Modern Math}
\usepackage{enumitem,needspace,xurl,hyperref}
\setlength{\parindent}{0pt}
\setlength{\parskip}{5pt}
\setlength{\emergencystretch}{3em}
\newcommand{\N}{\mathbb{N}}
\newcommand{\Z}{\mathbb{Z}}
\newcommand{\Q}{\mathbb{Q}}
\newcommand{\R}{\mathbb{R}}
\newcommand{\C}{\mathbb{C}}
\newcommand{\F}{\mathbb{F}}
\newcommand{\A}{\mathbb{A}}
\newcommand{\PP}{\mathbb{P}}
\newcommand{\E}{\mathbb{E}}
\newcommand{\eps}{\varepsilon}
\newcommand{\dd}{\,\mathrm{d}}
\begin{document}
\section*{UnsolvedMath ID: 30005963}
\textbf{Source problem code:} catalogue entry
\subsection{Definitions and mathematical statement}
Let \(\mathcal S\) denote the class called \emph{strict Calabi--Yau threefolds} (SCY3) in the supplied source. The record does not define this class or specify the base field; in particular, its intended smoothness, projectivity, canonical-bundle, fundamental-group, and cohomological hypotheses require confirmation from the cited discussion.

To make the numerical notation explicit, conditionally on the intended setting being smooth projective threefolds over \(\C\), write
\[
N^1(X)_{\Q}=\bigl(\operatorname{Div}(X)/{\equiv_{\mathrm{num}}}\bigr)\otimes_{\Z}\Q,
\qquad
N^1(X)_{\R}=N^1(X)_{\Q}\otimes_{\Q}\R.
\]
Let \(\operatorname{Amp}(X)\subseteq N^1(X)_{\R}\) be the cone of ample real divisor classes and set
\[
\operatorname{Amp}(X)_{\Q}=\operatorname{Amp}(X)\cap N^1(X)_{\Q}.
\]
With \(c_2(X)=c_2(T_X)\), define
\[
c_2(X)^\perp
=\{\alpha\in N^1(X)_{\R}:c_2(X)\cdot\alpha=0\},
\qquad
c_2(X)\cdot[D]=\int_X c_2(T_X)\smile c_1(\mathcal O_X(D)).
\]
Here \(\operatorname{Aut}(X)\) denotes the group of biregular automorphisms.

\textbf{Question 1.} Is the following assertion true?
\[
\forall X\in\mathcal S,\qquad
|\operatorname{Aut}(X)|=\infty
\ \Longrightarrow\
\exists\alpha\in\operatorname{Amp}(X)_{\Q}\setminus\{0\}
\quad c_2(X)\cdot\alpha=0.
\]
The word ``ample'' and the exclusion of zero are retained literally from the supplied formulations. Whether the source intends this cone or a different cone requires verification.

\textbf{Question 2.} Does there exist \(X\in\mathcal S\), other than the source's examples \(X_3\) and \(X_7\), with an automorphism \(f\in\operatorname{Aut}(X)\) that is primitive and has positive entropy?

For a complex projective interpretation, positive entropy means \(h_{\mathrm{top}}(f)>0\). A usual interpretation of \emph{primitive} is that there are no normal projective variety \(Y\), dominant rational map \(\pi:X\dashrightarrow Y\), and birational self-map \(g:Y\dashrightarrow Y\) such that
\[
0<\dim Y<3,
\qquad
\pi\circ f=g\circ\pi.
\]
This convention is conditional: the supplied record does not specify the intended definition of primitive, including any requirements on the fibres of \(\pi\). Likewise, it supplies neither constructions of \(X_3,X_7\) nor the equivalence relation intended by ``other than.'' Thus the conventional reading \(X\not\simeq X_3,X_7\) cannot be made fully explicit from this record alone. The original excerpt makes Question 2 existential; it does not ask whether every threefold satisfying Question 1's hypothesis admits such an automorphism.
\subsection{Short English statement}
The first question asks whether every strict Calabi--Yau threefold with infinitely many automorphisms has a nonzero rational ample divisor class pairing trivially with its second Chern class. The second asks for another strict Calabi--Yau threefold with a primitive automorphism of positive entropy, beyond the examples labelled \(X_3\) and \(X_7\). Definitions of the class of threefolds, primitivity, and the excluded examples require source confirmation.
\subsection{Sources}
\begin{itemize}
\item Local UnsolvedMath registry: supplied complete record for permanent ID \(30005963\), source problem code \texttt{catalogue entry}, titled \emph{Second-Chern Orthogonality and Calabi--Yau Automorphisms}, in the Oberwolfach Reports Open Problems collection. No individual problem proposer is supplied; the proposed year is recorded as 2024.
\item Supplied source citation: \emph{Algebraic Geometry: Wall Crossing and Moduli Spaces, Varieties and Derived Categories} (2025), Oberwolfach Reports, pp.~1787--1854, DOI \texttt{10.4171/owr/2024/32}. Original OWR report: \url{https://doi.org/10.4171/owr/2024/32}. The supplied original excerpt carries page number 1823.
\item Supplied literature source labelled ``Oguiso primary report'': \url{https://arxiv.org/abs/2407.17297}. No further bibliographic details are supplied.
\item Literature assessment dated exactly \texttt{2026-08-22}: the supplied curation labels both questions open and says that they remain open in the cited primary discussion. This is supplied dated curation, not an independently verified assessment. The background explicitly requests independent verification of time-sensitive status and cited claims.
\item The registry labels the statement \texttt{corrected\_verified} and reports that it was reconstructed and checked against the official Oberwolfach report. This records the supplied verification claim; no independent source consultation or mathematical review is asserted here.
\end{itemize}
\end{document}
